\documentclass[12pt]{article}

% Core math
\usepackage{amsmath, amssymb, amsfonts, amsthm}

% Page layout (load early so other packages see geometry)
\usepackage[margin=1in]{geometry}

% Graphics and color
\usepackage{graphicx}
\usepackage{xcolor}

% Diagrams
\usepackage{tikz}
\usepackage{tikz-cd}
\usetikzlibrary{arrows.meta, calc, positioning}

% Tables
\usepackage{booktabs}

% Page hooks for the GrokRxiv sidebar
\usepackage{everypage}

% Unicode characters used inside verbatim Lean snippets.
\usepackage{newunicodechar}
\newunicodechar{ℂ}{\ensuremath{\mathbb{C}}}
\newunicodechar{ℕ}{\ensuremath{\mathbb{N}}}
\newunicodechar{ℝ}{\ensuremath{\mathbb{R}}}
\newunicodechar{ℤ}{\ensuremath{\mathbb{Z}}}
\newunicodechar{ℚ}{\ensuremath{\mathbb{Q}}}
\newunicodechar{∀}{\ensuremath{\forall}}
\newunicodechar{∃}{\ensuremath{\exists}}
\newunicodechar{≠}{\ensuremath{\neq}}
\newunicodechar{≤}{\ensuremath{\leq}}
\newunicodechar{≥}{\ensuremath{\geq}}
\newunicodechar{π}{\ensuremath{\pi}}
\newunicodechar{ζ}{\ensuremath{\zeta}}
\newunicodechar{∧}{\ensuremath{\wedge}}
\newunicodechar{∨}{\ensuremath{\vee}}

% References (hyperref must be loaded late; cleveref after hyperref)
\usepackage{hyperref}
\usepackage{cleveref}

% Theorems
\newtheorem{theorem}{Theorem}[section]
\newtheorem{proposition}[theorem]{Proposition}
\newtheorem{lemma}[theorem]{Lemma}
\newtheorem{corollary}[theorem]{Corollary}
\newtheorem{conjecture}[theorem]{Conjecture}
\theoremstyle{definition}
\newtheorem{definition}[theorem]{Definition}
\newtheorem{problem}[theorem]{Problem}
\newtheorem{example}[theorem]{Example}
\theoremstyle{remark}
\newtheorem{remark}[theorem]{Remark}

% GrokRxiv sidebar
\definecolor{grokgray}{RGB}{110,110,110}
\AddEverypageHook{%
  \ifnum\value{page}=1
    \begin{tikzpicture}[remember picture, overlay]
      \node[
        rotate=90,
        anchor=south,
        font=\Large\sffamily\bfseries\color{grokgray},
        inner sep=0pt
      ] at ([xshift=38pt, yshift=0.52\paperheight]current page.south west)
      {GrokRxiv:2026.05.langlands-analytic-number-theory\quad
       [\,math.NT\,]\quad
       04 May 2026};
    \end{tikzpicture}
  \fi
}

% Custom macros
\newcommand{\C}{\mathbb{C}}
\newcommand{\R}{\mathbb{R}}
\newcommand{\Q}{\mathbb{Q}}
\newcommand{\Z}{\mathbb{Z}}
\newcommand{\N}{\mathbb{N}}
\newcommand{\F}{\mathbb{F}}
\newcommand{\HoTT}{\mathsf{HoTT}}
\newcommand{\Type}{\mathcal{U}}
\newcommand{\Prop}{\mathsf{Prop}}
\newcommand{\Set}{\mathsf{Set}}
\newcommand{\id}{\mathsf{id}}
\newcommand{\refl}{\mathsf{refl}}
\newcommand{\ua}{\mathsf{ua}}
\newcommand{\Path}{\mathsf{Path}}
\newcommand{\Cc}{\C_{\mathsf{c}}}
\newcommand{\Rc}{\R_{\mathsf{c}}}
\newcommand{\HIIT}{\mathsf{HIIT}}
\newcommand{\op}{\mathsf{op}}
\newcommand{\Aut}{\mathsf{Aut}}
\newcommand{\Hom}{\mathsf{Hom}}
\newcommand{\Map}{\mathsf{Map}}
\newcommand{\Eq}{\mathsf{Eq}}
\newcommand{\IsEquiv}{\mathsf{IsEquiv}}
\newcommand{\nCat}[1]{(\infty,#1)\text{-}\mathsf{Cat}}
\newcommand{\Tw}{\mathcal{T}}
\newcommand{\holo}{\mathsf{Holo}}
% Cohesive HoTT modalities.
\newcommand{\esh}{\mathsf{S}}      % shape modality
\newcommand{\flatm}{\mathsf{b}}    % flat modality (avoid clash with \flat)
\newcommand{\sharpm}{\mathsf{S^\sharp}} % sharp modality
% (We intentionally retain the standard \zeta from amssymb.)

\title{Toward HoTT-Native Analytic Number Theory:\\
Riemann Zeta, Langlands, and the $\zeta(s)=0$ Question}

\author{Matthew Long \\
\textit{The YonedaAI Collaboration} \\
\textit{YonedaAI Research Collective} \\
Chicago, IL \\
\texttt{research@yonedaai.com} $\cdot$ \url{https://yonedaai.com}}

\date{4 May 2026}

\begin{document}
\maketitle

\begin{abstract}
We address what the synthesis of our prior series of papers identified as the
\emph{principal open problem} in homotopy type theory's interface with
contemporary mathematics: the absence of a HoTT-native formulation of analytic
number theory. Algebraic number theory and arithmetic geometry, through the
lens of $(\infty,1)$-toposes, étale cohomology, and condensed mathematics,
admit clean translations into univalent foundations. By contrast, the Riemann
zeta function $\zeta(s)$, Dirichlet $L$-functions, automorphic forms, and the
Langlands programme have not been substantially reformulated in HoTT. The
statement $\zeta(s)=0$, viewed as a HoTT-native proposition with $\zeta$ a
HoTT-native object, remains unrealised. This paper does not solve that problem;
it formulates it precisely and offers a research roadmap.

We proceed in seven steps. First, we trace the prerequisite chain that any
HoTT-native definition of $\zeta$ must respect: HoTT-native Cauchy reals
$\Rc$, Cauchy complex numbers $\Cc$ via univalent algebraic closure,
holomorphic functions as a synthetic notion in cohesive HoTT, and Dirichlet
series as analytic objects internal to a univalent universe. Second, we
propose three candidate HoTT-native definitions of $\zeta$ — as a higher
inductive--inductive type, as the analytic limit of an Euler product, and as
the unique solution to a meromorphic continuation universal property — and
analyse their tradeoffs. Third, we situate this work against the geometric
Langlands proof of Gaitsgory--Raskin--Rozenblyum (2024), Clausen--Scholze
condensed and analytic-stack mathematics (2018--2025), and the
Loeffler--Stoll Lean 4 / Mathlib formalization of the Riemann zeta function
and $L$-functions (2025, arXiv:2503.00959). Fourth, we identify six concrete
sub-problems whose resolution would constitute a HoTT-native proof of
basic facts about $\zeta(s)$, including the Euler product, the functional
equation, and the formal statement of the Riemann hypothesis. We do not claim
the Riemann hypothesis, nor that HoTT will prove it; we claim that the
question of stating it inside HoTT is itself a non-trivial research programme,
and we take the first concrete steps toward that programme.
\end{abstract}

\tableofcontents

\section{Introduction}\label{sec:intro}

\subsection{The principal open problem}

In the synthesis paper concluding our prior series~\cite{Long2026Synthesis},
which builds on Riemann's foundational
1859~paper~\cite{RiemannHypothesisOriginal},
we surveyed six topics where homotopy type theory (HoTT) interfaces — or
fails to interface — with mainstream twenty-first century mathematics.
Topics 1, 3, 4, 5, and 6 (cohesive geometry, $\infty$-toposes, derived
algebraic geometry, motives, and condensed algebra) admit varying degrees of
HoTT formulation. Topic 2, \emph{analytic number theory}, does not. We wrote:
\begin{quote}
\itshape
Algebraic number theory and parts of arithmetic geometry — schemes, sheaves,
étale cohomology — admit $(\infty,1)$-topos and condensed mathematics
treatments that connect cleanly to univalent foundations. But $\zeta(s)=0$
as a HoTT-native statement, with $\zeta$ a HoTT-native object, remains
unrealised. We frame this as the principal open question for the synthesis.
\end{quote}
This paper develops that question.

\subsection{Why \texorpdfstring{$\zeta(s)=0$}{zeta(s)=0} is hard in HoTT}

A HoTT-native statement of $\zeta(s)=0$ requires four ingredients, each of
which is non-trivial:
\begin{enumerate}
\item A HoTT-native object $\Cc:\Type$ of complex numbers, complete with the
notion of meromorphicity. Constructive Cauchy reals~\cite{HoTTBook,Booij2020}
exist, but a univalent algebraic closure that is well-behaved with respect
to univalence requires care: classical algebraic closures are
unique-up-to-isomorphism, while a HoTT-native one is unique up-to-equivalence
in the sense of univalence~\cite{Voevodsky2014}.
\item A HoTT-native notion of \emph{holomorphic function} $\holo:(\Cc \to \Cc) \to \Type$.
The classical definition uses a metric limit (Newton quotient) which is
constructive but not synthetic. A more native approach uses cohesive HoTT
\cite{Shulman2018Cohesive}: holomorphicity becomes a modal predicate.
\item A HoTT-native definition of $\zeta:\Cc \setminus \{1\} \to \Cc$ that
agrees on $\mathrm{Re}(s)>1$ with $\sum_{n\geq 1} n^{-s}$ and analytically
continues elsewhere. This requires either a higher inductive-inductive
construction (HIIT) or an internal universal property of meromorphic
extensions.
\item A HoTT-native interpretation of the proposition
\[
\Pi(s:\Cc).\, \zeta(s) = 0 \to \neg\,\mathsf{trivialZero}(s) \to \mathrm{Re}(s) = 1/2.
\]
Even getting to this statement is a research project.
\end{enumerate}

By contrast, in classical foundations Loeffler and Stoll
\cite{LoefflerStoll2025} have just (March 2025) formalized in Lean 4 /
Mathlib the entire infrastructure: $\zeta$, the analytic continuation, the
functional equation, the Basel problem, non-vanishing on $\mathrm{Re}(s) \geq 1$,
Dirichlet's theorem on primes in arithmetic progressions, and the formal
statement of the Riemann hypothesis. Their work runs to roughly 3300 lines
of Lean for the analytic continuation alone. We have nothing comparable in
any HoTT proof assistant.

\subsection{What this paper does (and does not do)}

This paper:
\begin{itemize}
\item Identifies and motivates the \emph{prerequisite chain} for a HoTT-native
$\zeta$. (\Cref{sec:prereq})
\item Proposes \emph{three candidate} HoTT-native definitions of $\zeta$,
analyses their tradeoffs, and conjecturally proves their pairwise
equivalence. (\Cref{sec:zeta})
\item Places the picture in the broader geometric Langlands landscape via
$(\infty,1)$-toposes and Lurie's higher topos theory~\cite{Lurie2009HTT,
Gaitsgory2024GLC}. (\Cref{sec:langlands})
\item Compares HoTT and Clausen--Scholze condensed mathematics, identifying
where they overlap and where the bridge is missing. (\Cref{sec:condensed})
\item Compares with the Loeffler--Stoll 2025 Lean formalization, both as a
benchmark and as a \emph{translation target}. (\Cref{sec:loefflerstoll})
\item Provides a six-sub-problem roadmap that, if solved, would yield a
HoTT-native proof that $\zeta(2) = \pi^2/6$. (\Cref{sec:roadmap})
\item Re-states the Riemann hypothesis as a HoTT proposition and discusses
its modal-logical structure. (\Cref{sec:rh})
\end{itemize}

This paper does \emph{not}:
\begin{itemize}
\item Prove the Riemann hypothesis.
\item Provide a complete HoTT formalization of $\zeta$.
\item Claim that HoTT methods will, in principle, prove RH; we are agnostic
on this question.
\end{itemize}

The companion code repository contains a Haskell implementation of partial
zeta sums, Dirichlet series operations, the formal Euler product identity
(modulo a convergence axiom), and a Lean 4 sketch indexed against
Loeffler--Stoll's Mathlib code.

\subsection{Algebraic vs.\ analytic number theory}

A clarifying distinction. \emph{Algebraic number theory} studies
$\mathcal{O}_K$, the ring of integers of a number field $K/\Q$, its ideals,
class group, units, and reductions modulo primes. The objects are
discrete-finite or, at worst, countable; their morphisms are algebraic; and
the topology is the Zariski or étale one. HoTT formalisation is
straightforward in principle: rings are sets with operations satisfying
identities, and ideals are subsingletons of the underlying set.

\emph{Analytic number theory} studies $\zeta(s)$, Dirichlet $L$-functions
$L(s,\chi)$, and automorphic $L$-functions $L(s,\pi)$, where $\pi$ is an
automorphic representation of $GL(n,\mathbb{A}_\Q)$. The objects are
holomorphic / meromorphic functions of complex variables; morphisms are
analytic transformations; and the topology is the Euclidean (or condensed)
topology on $\C$. HoTT formalisation is presently absent.

The Langlands programme~\cite{Langlands1970,GelbartLanglands1979} unifies
these worlds: it asserts a precise correspondence between automorphic
representations and Galois representations, transforming questions about
zeros of $L$-functions into questions about Galois cohomology. Geometric
Langlands~\cite{Gaitsgory2024GLC} is the function-field analogue and was
proven in 2024 in five papers using $(\infty,1)$-categorical methods
throughout.

\subsection{Outline}

\Cref{sec:prereq} traces the prerequisite chain.
\Cref{sec:zeta} defines $\zeta$ as a candidate HoTT object.
\Cref{sec:langlands} treats geometric Langlands in $(\infty,1)$-topoi.
\Cref{sec:condensed} compares condensed mathematics and HoTT.
\Cref{sec:loefflerstoll} compares with Loeffler--Stoll Lean.
\Cref{sec:roadmap} gives the six-sub-problem roadmap.
\Cref{sec:rh} discusses RH as a HoTT statement.
\Cref{sec:conclusion} concludes.

%---------------------------------------------------------------
\section{The Prerequisite Chain}\label{sec:prereq}

We trace, with care, the dependency chain culminating in a HoTT-native $\zeta$.

\subsection{HoTT-native real numbers \texorpdfstring{$\Rc$}{Rc}}

\begin{definition}[Cauchy reals, after the HoTT Book~\cite{HoTTBook}]\label{def:Rc}
The type $\Rc$ of \emph{Cauchy reals} is the higher inductive-inductive type
defined simultaneously with a relation $\sim_\epsilon : \Rc \to \Rc \to \Type$
indexed by $\epsilon : \Q^+$, by the constructors:
\begin{align}
\mathsf{rat} &: \Q \to \Rc \\
\mathsf{lim} &: \Pi(x:\Q^+ \to \Rc).\, \mathsf{IsCauchy}(x) \to \Rc \\
\mathsf{eq} &: \Pi(u,v:\Rc).\, (\Pi(\epsilon:\Q^+).\, u \sim_\epsilon v) \to u = v
\end{align}
together with the path constructors for the closeness relation
$\sim_\epsilon$ \cite[Definition 11.3.2]{HoTTBook}.
\end{definition}

\begin{remark}
$\Rc$ is a set in the sense of HoTT (its identity types are propositions),
because of the explicit equality constructor $\mathsf{eq}$. This is in
contrast to a non-Cauchy approach where $\R$ might inherit higher structure
from its construction.
\end{remark}

\begin{theorem}[\cite{HoTTBook}, Theorem 11.3.32]\label{thm:RcArchimedean}
$\Rc$ is an Archimedean ordered field with decidable order $<$ on rationals
but undecidable equality on $\Rc$ in general.
\end{theorem}

\subsection{HoTT-native complex numbers \texorpdfstring{$\Cc$}{Cc}}

There are at least three approaches:

\paragraph{Pairs.}
$\Cc := \Rc \times \Rc$ with multiplication
$(a,b)(c,d) := (ac-bd, ad+bc)$. Simple, but does not capture the
\emph{universal property} that $\Cc$ should be an algebraic closure of $\Rc$.

\paragraph{Algebraic closure of $\Rc$.}
\begin{definition}[Univalent algebraic closure]\label{def:UAC}
A type $\overline{\Rc}$ together with an embedding $\iota:\Rc \to \overline{\Rc}$
is a \emph{univalent algebraic closure} of $\Rc$ if it is the propositional
truncation of the type
\[
  \sum_{F:\mathsf{Field}}\, \sum_{\iota:\Rc \to F} \mathsf{IsAlgClosure}(\Rc,F,\iota),
\]
where $\mathsf{IsAlgClosure}$ is the conjunction of (i) every non-zero
polynomial in $F[X]$ has a root, and (ii) every element of $F$ is algebraic
over $\Rc$ via $\iota$.
\end{definition}

\begin{remark}[Plain-language gloss]\label{rem:UACplain}
In essence, \Cref{def:UAC} defines the algebraic closure not as one
specific construction, but as the abstract object satisfying the universal
property of being an algebraically closed field containing $\Rc$. Univalence
ensures that any two such objects are propositionally equal as types in
$\mathcal{U}$. The propositional truncation $\|\cdot\|_{-1}$ at the top of
the $\Sigma$-type collapses the choice-of-witness data to a proposition,
reflecting the classical dictum that the algebraic closure is unique up to
non-canonical isomorphism.
\end{remark}

\begin{proposition}\label{prop:UACUnique}
Univalent algebraic closures of $\Rc$ are unique up to equivalence: if
$(\overline{\Rc}_1,\iota_1)$ and $(\overline{\Rc}_2,\iota_2)$ are two such,
there is an equivalence $e:\overline{\Rc}_1 \simeq \overline{\Rc}_2$ with
$e \circ \iota_1 = \iota_2$, and by univalence
$\overline{\Rc}_1 = \overline{\Rc}_2$.
\end{proposition}

\begin{proof}[Proof sketch]
Classical algebraic-closure uniqueness uses Zorn's lemma. In HoTT, we use a
constructive variant: one shows that any two algebraic closures are
isomorphic via an inductive construction, and that the choice of isomorphism
is contractible up to the action of $\mathsf{Gal}(\overline{\Rc}/\Rc)$. The
propositional truncation in \Cref{def:UAC} ensures the existential is a
proposition; uniqueness follows. Full proof requires choice or constructive
algebraic-closure arguments~\cite{Mines1988}.
\end{proof}

\paragraph{Quotient construction.}
$\Cc := \Rc[X]/(X^2 + 1)$. Concretely realisable via Schwartz / set-quotient.

In what follows, we treat $\Cc$ axiomatically: a propositional Univalent
Algebraic Closure of $\Rc$, equipped with a continuous structure inherited
from $\Rc \times \Rc$.

\subsection{Holomorphic functions in HoTT}

The classical definition of holomorphicity is:
\[
f:\C \to \C \text{ is holomorphic at } z \quad\text{iff}\quad
\lim_{h \to 0} \frac{f(z+h) - f(z)}{h} \text{ exists.}
\]
This translates straightforwardly to $\Rc$- and $\Cc$-limits, but is not
HoTT-native in the sense of using only $\Type$, $=$, and the
universe-polymorphic constructors.

\begin{definition}[Constructive holomorphicity]\label{def:holo}
$f:\Cc \to \Cc$ is \emph{holomorphic} at $z:\Cc$ if there exists
$f'(z):\Cc$ such that
\[
\Pi(\epsilon:\Q^+).\,
\Sigma(\delta:\Q^+).\,
\Pi(h:\Cc).\, 0 < |h| < \delta \to
\left|\,\frac{f(z+h)-f(z)}{h} - f'(z)\,\right| < \epsilon.
\]
\end{definition}

\begin{remark}[Synthetic alternative]\label{rem:synth}
In \emph{cohesive HoTT}~\cite{Shulman2018Cohesive}, the universe carries
shape and flat modalities $\esh, \flat$. Holomorphicity becomes a synthetic
predicate: $f$ is holomorphic iff $f$ commutes with the differential cohesion
operator. This is closer to ``HoTT-native'' but requires fixing a model
where $\esh$ is interpreted as the geometric $\C^\infty$-shape.
\end{remark}

\subsection{Dirichlet series}

\begin{definition}[Dirichlet series]\label{def:Dirichlet}
A \emph{Dirichlet series} is a function $a:\N_{>0} \to \Cc$ together with the
formal expression
\[
D_a(s) := \sum_{n=1}^{\infty} \frac{a(n)}{n^s}, \qquad s:\Cc.
\]
The \emph{abscissa of absolute convergence} is
$\sigma_a := \inf\{\sigma:\Rc \mid \sum_n |a(n)| n^{-\sigma} \text{ converges}\}$.
\end{definition}

\begin{proposition}[HoTT-native; sketched]\label{prop:DirichletAlgebra}
The set of Dirichlet series with $\sigma_a < +\infty$ forms a commutative
ring under termwise addition and Dirichlet convolution
$(a*b)(n) := \sum_{d \mid n} a(d) b(n/d)$.
\end{proposition}

\begin{proof}[Proof sketch]
Closed-under-addition is clear. For convolution, the abscissa of
$a*b$ is bounded by $\max(\sigma_a, \sigma_b) + \log_2(\text{divisor sum})$,
which is finite when both $\sigma_a, \sigma_b$ are. The full proof in HoTT
uses the absolute-convergence variant of Fubini for $\Rc$-valued double sums
\cite{Booij2020}.
\end{proof}

\subsection{Putting it together}

\begin{figure}[htbp]
\centering
\begin{tikzcd}[column sep=tiny, row sep=large]
\Q \arrow[r, hook] & \Rc \arrow[r, "\mathsf{UAC}"] & \Cc
  \arrow[r, "\holo"] & \mathsf{Hol}(\Cc, \Cc)
  \arrow[r, "\sum\frac{a(n)}{n^s}"]
  & \mathsf{Dir}(\Cc) \arrow[r, "\text{cont.}"] & \zeta
\end{tikzcd}
\caption{The prerequisite chain. Each arrow represents a non-trivial
HoTT-native construction. The final arrow ($\mathsf{Dir} \to \zeta$) is the
analytic continuation, and is the place at which the chain currently breaks
in pure HoTT.}
\label{fig:chain}
\end{figure}

The chain in \Cref{fig:chain} terminates at $\zeta$, but the rightmost arrow
— analytic continuation from a Dirichlet series convergent on $\mathrm{Re}(s)>1$
to a meromorphic function on $\Cc \setminus \{1\}$ — has no canonical
HoTT-native realisation. This is the central technical gap of the present
paper.

\subsection{Detailed account of analytic-continuation obstacles}\label{subsec:obstacles}

To clarify why the rightmost arrow is hard, we list the obstacles in
ascending order of severity.

\begin{enumerate}
\item \textbf{Unique factorisation of meromorphic extension.}
In classical complex analysis, the identity theorem (\Cref{thm:Identity})
guarantees uniqueness of analytic continuation. In HoTT, ``connected open''
needs a precise definition, and the proof that the agreement locus is
clopen requires constructive analogues of Bolzano--Weierstrass.

\item \textbf{Existence of meromorphic extension.}
Classical proofs of analytic continuation of $\zeta$ use one of:
\begin{itemize}
\item \emph{Hurwitz zeta} method: define $\zeta(s,a) = \sum_{n \geq 0}(n+a)^{-s}$
and continue via the integral representation. This requires Mellin transform.
\item \emph{Theta-function} method: $\theta(t) := \sum_{n \in \Z} e^{-\pi n^2 t}$
satisfies modular invariance, and Mellin transforming gives $\zeta$.
\item \emph{Contour integral} method: integrate $\frac{x^{s-1}}{e^x - 1}$
around a Hankel contour.
\end{itemize}
Each requires a substantial fragment of HoTT-native analysis: integration
theory, modular transformation laws, residue calculus.

\item \textbf{Computation in HoTT.}
Even given the existence and uniqueness, \emph{computing} $\zeta$ at, say,
$s = -1$ to obtain $\zeta(-1) = -1/12$ requires either (i) symbolic
manipulation of the analytic continuation, or (ii) numerical methods. HoTT
proof assistants like Cubical Agda do not currently support efficient
arbitrary-precision real arithmetic.

\item \textbf{Naturality.}
Classical proofs treat $\zeta$ as one specific function. A fully-categorified
HoTT-native account would treat $\zeta$ as part of a family — Dirichlet
$L$-functions, Hurwitz zeta, Hecke $L$-functions, automorphic
$L$-functions — with naturality with respect to characters and lifts. This
expands the existence obstacle by a constant factor of the size of the family.
\end{enumerate}

\subsection{HoTT-native Cauchy reals: explicit example}\label{subsec:CauchyEx}

To make the discussion concrete, we sketch how $\sqrt{2}:\Rc$ is a
HoTT-native object.

\begin{example}[$\sqrt{2}$ as Cauchy real]\label{ex:sqrt2}
Define $a:\Q^+ \to \Q$ by Newton iteration: $a(\epsilon) = $ the $n$-th
Newton iterate of $x_{n+1} := (x_n + 2/x_n)/2$, where $n$ is large enough
that $|a(\epsilon)^2 - 2| < \epsilon$. Then $a$ is a Cauchy modulus and
$\mathsf{lim}(a, \mathsf{IsCauchy}_a):\Rc$ is $\sqrt{2}$. Equality
$(\sqrt{2})^2 = 2$ in $\Rc$ follows from the equality constructor
$\mathsf{eq}$.
\end{example}

\begin{example}[$e$ as Cauchy real]\label{ex:euler}
$e := \mathsf{lim}((\epsilon \mapsto \sum_{k=0}^{N(\epsilon)} 1/k!),
\mathsf{IsCauchy}_e):\Rc$, where $N(\epsilon)$ is chosen so the tail
$\sum_{k>N} 1/k! < \epsilon$. By Stirling, $N(\epsilon) = O(\log(1/\epsilon)/
\log\log(1/\epsilon))$.
\end{example}

\begin{example}[$\pi$ as Cauchy real]\label{ex:pi}
We define
\[
\pi \;:=\; \mathsf{lim}\!\bigl(
   \epsilon \mapsto \text{Machin formula at depth $N(\epsilon)$},\;
   \mathsf{IsCauchy}_\pi\bigr).
\]
Machin's formula $\pi/4 = 4\arctan(1/5) - \arctan(1/239)$ converges
geometrically with ratios $1/25$ and $1/239^2$, so
$N(\epsilon) = O(\log(1/\epsilon))$.
\end{example}

These three examples illustrate that named real-number constants in HoTT
require an algorithmic Cauchy modulus, not just an existential statement.

%---------------------------------------------------------------
\section{The Riemann Zeta Function as a HoTT Object}\label{sec:zeta}

We now propose three approaches to defining $\zeta$ inside HoTT.

\subsection{Approach 1: \texorpdfstring{$\zeta$}{zeta} as a higher inductive-inductive type}

Inspired by the construction of $\Rc$ as a HIIT, we sketch a
\emph{specification pattern} for what a HoTT-native $\zeta$ should
satisfy. We emphasise that what follows is a \emph{wish-list of
constructors}, not a self-contained HIIT definition; both consistency
and existence of such a HIIT are open questions.

First, define the convergent partial-sum function on the half-plane of
absolute convergence:
\begin{equation}\label{eq:zetaSeries}
\zeta_{\mathrm{series}} : \Pi(s:\Cc).\, \mathrm{Re}(s) > 1 \to \Cc, \qquad
\zeta_{\mathrm{series}}(s) := \sum_{n=1}^{\infty} n^{-s}.
\end{equation}
This map is well-defined because $\mathrm{Re}(s) > 1$ ensures absolute
convergence of the series in $\Cc$, and the limit is the (HoTT-native)
limit constructor of $\Rc$ applied component-wise.

\begin{definition}[Specification of zeta as HIIT, candidate]\label{def:ZetaHIIT}
A pair $(\zeta:\Cc \setminus \{1\} \to \Cc, P:\mathsf{Holo}(\zeta))$ is a
\emph{HIIT-realisation of zeta} if it is generated by:
\begin{align}
\zeta_{\mathrm{Re}>1} &: \Pi(s:\Cc).\, \mathrm{Re}(s) > 1 \to \Cc, \notag \\
\zeta_{\mathrm{Re}>1}(s) &= \zeta_{\mathrm{series}}(s) \quad
   \text{(equality with the series of \cref{eq:zetaSeries})} \\
\mathsf{cont} &: \Pi(s:\Cc).\, s \neq 1 \to \Cc \\
\mathsf{agree} &: \Pi(s:\Cc).\, \Pi(p:\mathrm{Re}(s) > 1).\,
   \mathsf{cont}(s, \mathrm{ne1}_p) = \zeta_{\mathrm{Re}>1}(s, p) \\
\mathsf{holo} &: \mathsf{Holo}(\mathsf{cont})
\end{align}
together with \emph{additional path-constructors} (intentionally left
schematic) enforcing the functional equation
$\zeta(1-s) = 2(2\pi)^{-s}\Gamma(s)\cos(\pi s/2)\zeta(s)$.

We use the language of a \texttt{definition} only for the specification;
we make no claim of consistency.
\end{definition}

\begin{remark}\label{rem:ZetaHIITissues}
This is a \emph{specification}, not a finished HIIT, for two reasons:
\begin{enumerate}
\item The $\mathsf{holo}$ constructor demands a holomorphic extension, but no
such extension is canonically given by the constructors themselves.
Consistency therefore reduces to a separate \emph{existence-of-extension}
lemma — which is exactly the analytic-continuation gap of \Cref{prob:4}.
\item The functional-equation path-constructors are intentionally schematic.
Their precise form would involve HoTT-native $\Gamma$, $\cos$, and the
modular-transformation identity for $\theta$, none of which are presently
formalised. We do not know what their definitive shape should be, and
indicate this rather than papering over it.
\end{enumerate}
The specification is therefore a \emph{target pattern}, not a finished
construction. The remaining two approaches (\Cref{def:ZetaEuler} and
\Cref{def:ZetaUP}) are more conservative.
\end{remark}

\subsection{Approach 2: \texorpdfstring{$\zeta$}{zeta} as analytic limit of Euler product}

Recall the Euler product (\cite{Euler1737}):
\[
\zeta(s) = \prod_{p \text{ prime}} \frac{1}{1 - p^{-s}}, \qquad \mathrm{Re}(s) > 1.
\]

\begin{definition}[Zeta via Euler product]\label{def:ZetaEuler}
On the half-plane $H_1 := \{s:\Cc \mid \mathrm{Re}(s) > 1\}$, define
\[
\zeta_E(s) := \prod_{p \in \mathbb{P}} (1 - p^{-s})^{-1}
\]
where $\mathbb{P}$ is the type of primes. Extend by analytic continuation
(separately proven) to $\Cc \setminus \{1\}$.
\end{definition}

\begin{proposition}[Euler product, modulo convergence]\label{prop:Euler}
On $H_1$, $\zeta_E(s) = \sum_{n=1}^{\infty} n^{-s}$.
\end{proposition}

\begin{proof}[Proof sketch in HoTT]
By unique factorisation in $\N_{>0}$ (a HoTT-formalisable theorem), each
$n \geq 1$ corresponds bijectively to a finite-support function $\mathbb{P}
\to \N$ via $n = \prod_p p^{e_p(n)}$. Expanding the Euler product:
\[
\prod_p (1 - p^{-s})^{-1}
  = \prod_p \sum_{k \geq 0} p^{-ks}
  = \sum_{(e_p)\in \N^{(\mathbb{P})}} \prod_p p^{-e_p s}
  = \sum_{n \geq 1} n^{-s}.
\]
The middle equality uses absolute convergence of the product on $H_1$, which
is a HoTT-statement requiring \Cref{prop:DirichletAlgebra}.
\end{proof}

\subsection{Approach 3: \texorpdfstring{$\zeta$}{zeta} via universal property of meromorphic continuation}

\begin{definition}[Meromorphic continuation, universal]\label{def:MeromCont}
Given a holomorphic $f:U \to \Cc$ on an open $U \subseteq \Cc$, a
\emph{meromorphic continuation} of $f$ to $V \supseteq U$ is a
meromorphic $\tilde f:V \to \Cc$ with $\tilde f|_U = f$, such that for any
other meromorphic continuation $g:V \to \Cc$, $\tilde f = g$ as meromorphic
functions. (Existence requires connectedness; uniqueness uses the identity
theorem.)
\end{definition}

\begin{theorem}[Identity theorem, HoTT version]\label{thm:Identity}
Suppose $f, g:V \to \Cc$ are holomorphic on a connected open $V \subseteq \Cc$,
and the type $\sum_{z:V} f(z) = g(z)$ has an accumulation point in $V$.
Then $\Pi(z:V).\, f(z) = g(z)$.
\end{theorem}

\begin{proof}[Proof sketch]
Standard: the locus of agreement is open (by power series), closed (by
continuity), and non-empty, hence equal to $V$ since $V$ is connected.
HoTT-native version: ``connected'' becomes ``$\|V\|_{-1}$ is a singleton''
and ``open'' is interpreted appropriately in cohesive HoTT.
\end{proof}

\begin{definition}[Zeta via universal property]\label{def:ZetaUP}
$\zeta:\Cc \setminus \{1\} \to \Cc$ is the unique (up to identity, by
\Cref{thm:Identity}) meromorphic continuation of $s \mapsto \sum n^{-s}$ from
$H_1$ to $\Cc \setminus \{1\}$ with a simple pole of residue $1$ at $s=1$.
\end{definition}

\subsection{Equivalence of the three definitions}

\begin{theorem}[Conjectural equivalence]\label{thm:ZetaEquiv}
Conditional on the existence of HoTT-native analytic continuation, the three
definitions \Cref{def:ZetaHIIT}, \Cref{def:ZetaEuler}, \Cref{def:ZetaUP}
are pairwise equivalent: there is a propositional equality between any two.
\end{theorem}

\begin{proof}[Proof sketch]
\Cref{def:ZetaEuler} agrees with the partial-sum definition on $H_1$ by
\Cref{prop:Euler}. \Cref{def:ZetaUP} agrees with the partial-sum definition
on $H_1$ by definition. By \Cref{thm:Identity}, any two meromorphic
continuations of the same function agree everywhere. \Cref{def:ZetaHIIT}
imposes both the partial-sum agreement and the holomorphicity constraint,
hence its result is identified with the universal-property zeta. The
catch: each step requires HoTT-native analytic continuation, which is
exactly the gap.
\end{proof}

\subsection{The functional equation}

\begin{theorem}[Functional equation, conjectured HoTT-native]\label{thm:FE}
For all $s : \Cc \setminus \{0,1\}$,
\[
\zeta(1-s) = 2 (2\pi)^{-s} \Gamma(s) \cos(\pi s / 2) \zeta(s).
\]
\end{theorem}

\begin{proof}[Strategy]
Riemann's original proof uses the theta-function identity
$\theta(t) = t^{-1/2}\theta(1/t)$ and the Mellin transform. Translating to
HoTT requires:
\begin{itemize}
\item HoTT-native theta function $\theta:\Rc^+ \to \Rc$, defined by
$\theta(t) := \sum_{n \in \Z} e^{-\pi n^2 t}$, plus the modular identity.
\item HoTT-native Mellin transform $\mathcal{M}:\holo(\Rc^+,\Rc) \to \holo(\Cc,\Cc)$.
\item HoTT-native contour-integral lemmas.
\end{itemize}
None of these are presently formalized in HoTT. They are all formalized in
classical Lean 4 / Mathlib (Loeffler--Stoll, 2025).
\end{proof}

\subsection{The critical strip}

\begin{definition}[Critical strip]
$S := \{s : \Cc \mid 0 < \mathrm{Re}(s) < 1\}$.
\end{definition}

\begin{definition}[Trivial zeros]
A zero $s_0$ of $\zeta$ is \emph{trivial} if $s_0 \in \{-2,-4,-6,\ldots\}$.
\end{definition}

\begin{conjecture}[Riemann hypothesis, HoTT statement]\label{conj:RH}
\[
\Pi(s:\Cc).\, \zeta(s) = 0 \to \neg\,\mathsf{trivialZero}(s) \to \mathrm{Re}(s) = 1/2.
\]
\end{conjecture}

We will analyse this statement modal-logically in \Cref{sec:rh}.

\subsection{Worked example: \texorpdfstring{$\zeta(2) = \pi^2/6$}{zeta(2) = pi squared over 6}}\label{subsec:Basel}

To illustrate what HoTT-native machinery is needed, we trace one of the
oldest results — the Basel problem.

\begin{theorem}[Basel problem]\label{thm:Basel}
$\zeta(2) = \pi^2/6$.
\end{theorem}

\begin{proof}[Sketch of HoTT-native proof, modulo missing infrastructure]
There are several classical approaches; we outline two.

\textbf{Approach (i): Fourier series of $x^2$.} Expand $x^2$ on $[-\pi,\pi]$
as a Fourier series:
\[
x^2 = \frac{\pi^2}{3} + 4\sum_{n=1}^{\infty} \frac{(-1)^n \cos(nx)}{n^2}.
\]
Setting $x = \pi$ yields $\pi^2 = \pi^2/3 + 4\zeta(2)$,
hence $\zeta(2) = \pi^2/6$. The HoTT-native version requires:
\begin{itemize}
\item HoTT-native Fourier series with pointwise convergence on smooth
functions.
\item HoTT-native trigonometric functions (definable as power series, hence
HIIT $\Rc$-valued).
\item Pointwise evaluation at the boundary requires Abel-style limits.
\end{itemize}

\textbf{Approach (ii): Euler's product expansion of $\sin$.} Use
\[
\frac{\sin(\pi z)}{\pi z} = \prod_{n=1}^{\infty}\left(1 - \frac{z^2}{n^2}\right).
\]
Equating Taylor coefficients of $z^2$ on both sides:
$-\pi^2/6 = -\sum 1/n^2$, so $\zeta(2) = \pi^2/6$. HoTT-native version
requires:
\begin{itemize}
\item Infinite-product convergence theory in HoTT.
\item Term-by-term Taylor expansion of the product (Abel-Mertens-style
manipulation).
\item Product-to-sum identity, valid on absolute convergence.
\end{itemize}

Either approach requires roughly the same prerequisite chain:
\Cref{prob:1}--\Cref{prob:3}, plus some elementary measure / convergence
theory. \Cref{prob:4} is not needed for this specific theorem because $s=2$
is in the half-plane of absolute convergence.
\end{proof}

\begin{remark}\label{rem:zetatwo}
\Cref{thm:Basel} is a \emph{lower-bound} HoTT goal: it does not require
analytic continuation or the functional equation. We propose it as the
\emph{minimum viable target} for HoTT-native analytic NT — the analogue of
``hello, world'' for our roadmap.
\end{remark}

%---------------------------------------------------------------
\section{Geometric Langlands in $(\infty,1)$-Topoi}\label{sec:langlands}

\subsection{Brief history}

The Langlands programme~\cite{Langlands1970} predicts a correspondence
\[
\bigl\{ \text{automorphic reps of } GL(n,\mathbb{A}_F) \bigr\}
\;\longleftrightarrow\;
\bigl\{ \text{$n$-dim Galois reps of } \mathrm{Gal}(\overline{F}/F) \bigr\}
\]
for a global field $F$. The number-field case ($F$ a number field) is
analytic in nature; the function-field case ($F = \F_q(X)$ for a curve $X$)
is geometric, hence the name \emph{geometric Langlands}.

Geometric Langlands has a categorical formulation due to Beilinson--Drinfeld
\cite{BeilinsonDrinfeld2004} and Frenkel--Gaitsgory: the conjecture is
the existence of an equivalence
\[
\mathsf{D}\text{-mod}(\mathsf{Bun}_G(X)) \,\simeq\, \mathsf{QCoh}(\mathsf{LocSys}_{G^\vee}(X))
\]
of $(\infty,1)$-categories, where $G$ is a reductive group, $\mathsf{Bun}_G$
its moduli stack of $G$-bundles, $G^\vee$ the Langlands dual, and
$\mathsf{LocSys}$ the de Rham moduli of local systems. This was
\textbf{proven} in 2024 by Gaitsgory, Raskin, Rozenblyum, Arinkin, Beraldo,
Chen, Cheng, Faergeman, Lin, Lysenko in five papers
\cite{Gaitsgory2024GLC}, awarded the 2025 Breakthrough Prize.

\subsection{\texorpdfstring{$(\infty,1)$}{(infty,1)}-categories vs.\ HoTT}

By Lurie's higher topos theory~\cite{Lurie2009HTT} and higher
algebra~\cite{LurieHA}, $(\infty,1)$-categories
admit a model in simplicial sets (quasi-categories). Cisinski et al.\
\cite{CisinskiKKNS2025} have begun a programme of \emph{type-theoretic
foundations of $(\infty,1)$-category theory} in which $(\infty,1)$-categories
are types satisfying a Segal condition.

By Shulman's theorem~\cite{Shulman2019InfTopos}, every $(\infty,1)$-topos
$\mathcal{X}$ admits HoTT as its internal language. Hence, in principle,
the Gaitsgory equivalence is statable in HoTT \emph{internally to a fixed
$(\infty,1)$-topos}.

\begin{remark}\label{rem:internal}
The catch is that the equivalence relates two \emph{different} topoi:
$\mathsf{D}\text{-mod}(\mathsf{Bun}_G)$ is internal to a derived geometric
topos; $\mathsf{QCoh}(\mathsf{LocSys}_{G^\vee})$ is internal to a different
derived geometric topos. A HoTT statement requires either an ambient
2-topos / very-large-universe setup (in the spirit of Riehl--Verity directed
type theory) or an external HoTT statement comparing two HoTT internal
languages.
\end{remark}

\subsection{Geometric vs.\ analytic Langlands}

The 2024 proof is \emph{geometric}: it lives in the world of moduli stacks
over $\F_q$-curves or, in characteristic 0, over $\C$-curves with the de
Rham stack (cf.~\cite{ClausenScholzeComplex} for the complex-analytic side). It says nothing directly about $\zeta$ or about analytic number
theory.

The \emph{analytic} Langlands programme (Etingof--Frenkel--Kazhdan
\cite{EFK2025AnalyticLanglands}) attempts a Hilbert-space variant: the
$L^2$-spectrum of certain operators on $\mathsf{Bun}_G(\C\text{-curve})$
should match a spectral side built from $\mathsf{LocSys}_{G^\vee}$. This is
the program closer to $\zeta(s) = 0$, but it is much less developed than
the geometric version.

\subsection{Implication for HoTT-native analytic NT}

If a HoTT-native analytic Langlands programme could be developed, then
$\zeta(s) = 0$ would translate to a statement about the spectrum of an
operator on $\mathsf{Bun}_G(\Q)$, which is at least \emph{syntactically} a
HoTT statement modulo standard moduli-stack constructions. We do not
develop this direction in detail; we flag it as a concrete research direction.

\subsection{Detailed example: \texorpdfstring{$GL(1)$}{GL(1)} Langlands and Hecke characters}\label{subsec:GL1}

The simplest case of Langlands is $G = GL(1)$, and even this case shows
where HoTT-native infrastructure is needed.

\begin{example}[$GL(1)$ over $\Q$]
Automorphic representations of $GL(1, \mathbb{A}_\Q)$ are continuous
characters $\chi: \mathbb{A}_\Q^\times / \Q^\times \to \C^\times$. By
class-field theory, these correspond to characters of $\mathrm{Gal}^{\mathrm{ab}}
(\overline{\Q}/\Q)$, i.e., characters of the idele class group. The
$L$-function attached to $\chi$ is
\[
L(s,\chi) = \prod_p \frac{1}{1 - \chi(p) p^{-s}}.
\]
For $\chi = 1$ trivial, this is $\zeta(s)$. The Langlands correspondence
identifies $L(s,\chi)$ with the Artin $L$-function of the corresponding
Galois character.
\end{example}

\begin{remark}[HoTT formulation of Hecke characters]
A HoTT formulation requires:
\begin{itemize}
\item HoTT-native ideles $\mathbb{A}_\Q^\times$ as a restricted product
over places. This is a HoTT-native colimit over a directed system of
finite-place subgroups; HoTT-native completion at each place requires HoTT
$p$-adic numbers~\cite{HoTTBook}.
\item HoTT-native continuous group homomorphisms.
\item HoTT-native $L$-function attached to character.
\end{itemize}
None of these are presently available, but each is plausibly a few thousand
lines of HoTT code.
\end{remark}

\subsection{Geometric Langlands as an internal HoTT statement}\label{subsec:GLCInternal}

We give a more precise version of the Gaitsgory equivalence as a HoTT
internal statement. Let $\mathcal{X}$ be the $(\infty,1)$-topos of derived
algebraic stacks over $\C$. By Shulman~\cite{Shulman2019InfTopos},
$\mathcal{X}$ has an internal language extending HoTT (after fixing a
universe-of-types issue).

\begin{conjecture}[Gaitsgory equivalence in HoTT, schematic]\label{conj:GLCHoTT}
In the internal language of $\mathcal{X}$, fix a smooth projective curve
$X:\mathcal{X}$ over $\C$ and a reductive group $G:\mathcal{X}$. Define:
\begin{align*}
\mathsf{Bun}_G(X) &:\Type \quad \text{(moduli stack of $G$-bundles)}, \\
\mathsf{LocSys}_{G^\vee}(X) &:\Type \quad \text{(moduli stack of $G^\vee$-local systems)}, \\
\mathsf{D}\text{-mod}(\mathsf{Bun}_G) &:\Type \quad \text{(category of D-modules)}, \\
\mathsf{QCoh}(\mathsf{LocSys}_{G^\vee}) &:\Type \quad \text{(category of quasi-coherent sheaves)}.
\end{align*}
There is a HoTT-internal equivalence of $\infty$-categories:
\[
\mathsf{D}\text{-mod}(\mathsf{Bun}_G(X)) \simeq \mathsf{QCoh}(\mathsf{LocSys}_{G^\vee}(X)).
\]
\end{conjecture}

This is conjectural in the sense that we have not verified all the
type-theoretic encodings; the underlying mathematical content is theorem
(Gaitsgory et al.\ 2024).

\begin{remark}\label{rem:univalence-glc}
Univalence enters when one asks: \emph{which} equivalence? The
Gaitsgory--Drinfeld equivalence comes equipped with a Hecke-eigensheaf
property; under univalence, this distinguishes one canonical equivalence up
to a contractible space of choices.
\end{remark}

\subsection{Physical interpretation: 4d \texorpdfstring{$\mathcal{N}=4$}{N=4} super-Yang--Mills}

Kapustin and Witten~\cite{KapustinWitten2007} interpreted geometric
Langlands as electric--magnetic duality (S-duality) of 4d
$\mathcal{N}=4$ super-Yang--Mills compactified on a Riemann surface.
This physical perspective suggests:
\begin{itemize}
\item A HoTT-native treatment of 4d $\mathcal{N}=4$ super-Yang--Mills via
synthetic differential cohesive HoTT~\cite{Schreiber2013}.
\item S-duality as an automorphism of the underlying type; eigenvalues of
S-duality giving the spectrum.
\item Connection to the synthesis paper's QFT formulation.
\end{itemize}
We flag this as a research direction; we do not develop it further here.

%---------------------------------------------------------------
\section{Condensed Mathematics and HoTT}\label{sec:condensed}

\subsection{Pyknotic / condensed sets}

Clausen and Scholze~\cite{Scholze2019Condensed,ClausenScholze2024Analytic}
defined the \emph{condensed sets} as sheaves on the site of profinite sets
with finite jointly-surjective covers. Pyknotic sets, due to Barwick--Haine,
are an essentially equivalent variant. The category $\mathsf{Cond}$ is a
topos, and $\mathsf{Cond}(\mathrm{Ab})$ — condensed abelian groups — has
\emph{much better homological properties} than topological abelian groups.

The crucial example: $\C$ becomes a \emph{condensed ring} which is
analytic-friendly in a way that $\C$-as-topological-ring is not. The
6-functor formalism of analytic stacks (Clausen--Scholze 2024) gives a
geometric foundation for analytic-number-theoretic objects.

\subsection{Why this matters for \texorpdfstring{$\zeta$}{zeta}}

The condensed approach gives a uniform setting in which:
\begin{itemize}
\item Smooth manifolds, complex-analytic spaces, schemes, formal schemes,
adic spaces, and rigid spaces all live as condensed objects.
\item Cohomology operations (six functors: $f^*, f_*, f_!, f^!, \otimes,
\mathrm{Hom}$) all exist with clean adjunction structure.
\item $\zeta$, viewed as a meromorphic function on $\C$, becomes an object
in $\mathsf{Cond}(\mathrm{Ab})$-modules over a suitable condensed analytic
stack.
\end{itemize}

\subsection{Bridging to HoTT}

\begin{problem}[Bridging condensed mathematics and HoTT]\label{prob:bridge}
Construct an $(\infty,1)$-topos $\mathcal{X}$ with an internal language
extending HoTT, in which:
\begin{itemize}
\item Condensed sets embed fully faithfully.
\item Solid abelian groups (in the sense of Clausen--Scholze) are an
internal type.
\item Holomorphic / meromorphic functions on $\C$ correspond to morphisms
of condensed analytic stacks.
\end{itemize}
\end{problem}

\Cref{prob:bridge} is open; preliminary work by Mahmoudvand--Riehl and others
has explored the syntactic side, but no complete bridge exists.

\subsection{Solid abelian groups in HoTT, sketch}

\begin{definition}[Solid abelian group, condensed]
A condensed abelian group $A$ is \emph{solid} if for every profinite set
$S$ and every continuous $S \to A$ which is null on the closure of zero,
the induced map factors through $A$ uniquely.
\end{definition}

A HoTT-native version would replace ``profinite set'' with a HoTT-internal
type (e.g., a limit of finite types in $\mathcal{U}$) and ``continuous''
with the appropriate cohesive-HoTT modality. We sketch a candidate:

\begin{definition}[HoTT-solid abelian group, candidate]\label{def:HoTTSolid}
Working in cohesive HoTT, an abelian group $A:\mathcal{U}$ is \emph{HoTT-solid}
if $\esh(A) \simeq A$ and the map $\flat A \to A$ is an equivalence.
\end{definition}

\begin{remark}
\Cref{def:HoTTSolid} is speculative; it has not been verified to match the
condensed definition. It is offered as a starting point for future work.
\end{remark}

\subsection{Six-functor formalism in HoTT}\label{subsec:sixfunctor}

The Clausen--Scholze 6-functor formalism for analytic
stacks~\cite{ClausenScholze2024SixFunctor,ClausenScholze2025AnalyticStacks}
is expressed in $(\infty,1)$-categorical language as follows. Given a category
$\mathcal{C}$ of geometric objects (analytic stacks) with two classes of
morphisms ``proper'' $\mathsf{P}$ and ``open'' $\mathsf{O}$, a 6-functor
formalism consists of:
\begin{itemize}
\item A symmetric monoidal $(\infty,1)$-category $\mathcal{D}(X)$ for each
$X \in \mathcal{C}$, of $\infty$-sheaves on $X$.
\item For each morphism $f: X \to Y$ in $\mathcal{C}$, an adjunction
$f^* \dashv f_*: \mathcal{D}(Y) \rightleftarrows \mathcal{D}(X)$.
\item For each $f$ in $\mathsf{P} \cup \mathsf{O}$, an additional pair
$f_! \dashv f^!: \mathcal{D}(X) \rightleftarrows \mathcal{D}(Y)$.
\item Coherent base change, projection formula, and proper / open
recollement.
\end{itemize}

\begin{problem}[6-functors in HoTT]\label{prob:sixfunc}
Express the 6-functor formalism in HoTT internal language. The natural
setting: $(\infty,1)$-topoi indexed over a directed type / $(\infty,2)$-base
\cite{Riehl2021Elements}.
\end{problem}

A successful resolution of \Cref{prob:sixfunc} would, indirectly, give a
HoTT framework for analytic objects (including $\zeta$ as a section of a
suitable line bundle).

\subsection{Solid modules over \texorpdfstring{$\Z$}{Z} and \texorpdfstring{$p$}{p}-adic interaction}\label{subsec:solidp}

In condensed mathematics, the category of solid abelian groups
$\mathrm{Solid}_{\Z}$ contains $p$-adic objects naturally:
\[
\Z_p \in \mathrm{Solid}_{\Z}, \qquad \Q_p \in \mathrm{Solid}_{\Z}, \qquad
\C_p \in \mathrm{Solid}_{\Z}.
\]
For $\Z_p$, the topology comes from the profinite structure
$\Z_p = \varprojlim_n \Z/p^n$. For $\C_p$, completing the algebraic closure
of $\Q_p$ uses metric completion plus algebraic closure plus completion
again — this is awkward classically and equally awkward in HoTT.

\begin{conjecture}[$p$-adic zeta function in HoTT]\label{conj:padicZeta}
A HoTT-native realisation of the Kubota--Leopoldt $p$-adic zeta function
$\zeta_p:\Z_p \setminus \{1\} \to \Q_p$ is plausibly easier than the
archimedean $\zeta$, because the $p$-adic case avoids analytic continuation
and uses interpolation of values $\zeta(1-n)$ for $n \geq 1$.
\end{conjecture}

\begin{remark}\label{rem:padicEasier}
\Cref{conj:padicZeta} is a research suggestion: \emph{start with the $p$-adic
case}. Loeffler--Stoll formalize the archimedean case; the $p$-adic case has
not yet been comprehensively formalized in any proof assistant. HoTT could
be the first.
\end{remark}

%---------------------------------------------------------------
\section{Comparison with Loeffler--Stoll Lean Formalization}\label{sec:loefflerstoll}

\subsection{What Loeffler--Stoll have}

The 2025 paper~\cite{LoefflerStoll2025} formalizes in Lean 4 / Mathlib:
\begin{enumerate}
\item The Riemann zeta function $\zeta$ via Hurwitz zeta and analytic
continuation, $\sim$ 3300 lines.
\item Dirichlet $L$-functions $L(s,\chi)$ for primitive characters $\chi$.
\item The Euler product $\zeta(s) = \prod_p (1-p^{-s})^{-1}$ on
$\mathrm{Re}(s)>1$.
\item Analytic continuation via Mellin / theta-function method.
\item The functional equation $\zeta(1-s) = 2(2\pi)^{-s}\Gamma(s)
\cos(\pi s/2)\zeta(s)$.
\item Basel: $\zeta(2) = \pi^2/6$.
\item Non-vanishing on $\mathrm{Re}(s) \geq 1$.
\item Dirichlet's theorem on primes in APs.
\item Formal statement of the Riemann hypothesis:
\begin{verbatim}
def RiemannHypothesis : Prop :=
  ∀ s : ℂ, riemannZeta s = 0 → ¬ trivialZero s → s.re = 1 / 2
\end{verbatim}
\end{enumerate}

This is roughly $10^4$ lines total across the relevant Mathlib files.

\subsection{What Lean / Mathlib provides that HoTT lacks}

\begin{itemize}
\item Classical $\C$ as a complete normed field, with Cauchy integrals.
\item Mathlib's \texttt{Complex.differentiable}, \texttt{HolomorphicAt},
\texttt{ContourIntegral}.
\item Mellin transforms, Gamma function, Riemann zeta function
(\texttt{ZetaFunction} in Mathlib).
\item All of measure theory, including Lebesgue integration on $\C$.
\end{itemize}

\subsection{What HoTT could offer}

\begin{itemize}
\item \emph{Univalence}: the multiple definitions of $\zeta$
(\Cref{def:ZetaHIIT}, \Cref{def:ZetaEuler}, \Cref{def:ZetaUP}) become
\emph{propositionally equal} via univalence, not merely equal up to
some isomorphism.
\item \emph{Higher inductive types}: the Riemann surface of $\zeta$ (or of
$L$-functions) is naturally a HIT, capturing branch-cut data
syntactically.
\item \emph{Internal $(\infty,1)$-topos language}: the Langlands functor and
its adjoints can be expressed in the same language as the underlying
analytic objects.
\end{itemize}

\subsection{Strategy: HoTT as classical-foundation transport}

A pragmatic approach: \emph{transport} the Loeffler--Stoll Lean
formalization to HoTT \emph{not} by re-doing all the analysis, but by
treating $\C$ axiomatically as a HoTT type with the same operations and
properties as Mathlib's $\C$. We sketch this in \Cref{sec:roadmap}.

\subsection{Code-level comparison: signature against signature}\label{subsec:LSsigs}

To make the comparison tangible, we present the key Lean signatures from
Loeffler--Stoll alongside conjectured HoTT analogues.

\paragraph{The zeta function itself.}
In Mathlib (Loeffler--Stoll):
\begin{verbatim}
noncomputable def riemannZeta : ℂ → ℂ
\end{verbatim}
In HoTT (proposed):
\[
\zeta : \Cc \setminus \{1\} \to \Cc \quad \text{(with explicit pole at $s=1$)}.
\]
Here Lean uses ``noncomputable'' to encode that $\zeta$ is defined by
non-effective analytic continuation. HoTT, lacking a built-in concept of
``noncomputable'', must encode this via a propositional truncation.

\paragraph{The functional equation.}
In Mathlib:
\begin{verbatim}
theorem riemannZeta_one_sub : ∀ (s : ℂ), s ≠ 0 → s ≠ 1 →
  riemannZeta (1 - s) = 2 * (2 * π) ^ (-s) * Gamma s
                      * cos (π * s / 2) * riemannZeta s
\end{verbatim}
HoTT analogue (proposed): exactly the corresponding $\Pi$-statement, modulo
the conditional $s \neq 0 \wedge s \neq 1$.

\paragraph{Non-vanishing on $\mathrm{Re}(s) \geq 1$.}
In Mathlib:
\begin{verbatim}
theorem riemannZeta_ne_zero_of_one_le_re :
  ∀ (s : ℂ), 1 ≤ s.re → s ≠ 1 → riemannZeta s ≠ 0
\end{verbatim}
HoTT analogue: same statement; the proof uses the Euler product (no
analytic continuation), so should be tractable HoTT-natively.

\paragraph{Riemann hypothesis.}
In Mathlib:
\begin{verbatim}
def RiemannHypothesis : Prop :=
  ∀ s : ℂ, riemannZeta s = 0 → ¬ trivialZero s → s.re = 1 / 2
\end{verbatim}
HoTT analogue: \Cref{conj:RH}.

\subsection{Quantitative comparison}\label{subsec:LSquant}

\begin{center}
\begin{tabular}{lcc}
\toprule
Theorem & Lean lines (Loeffler--Stoll) & Est.\ HoTT complexity \\
\midrule
$\zeta$ definition & 200 & $\sim 2$--$3\times$ \\
Analytic continuation & 3300 & $\sim 2$--$3\times$ \\
Functional equation & 800 & $\sim 2$--$3\times$ \\
Euler product & 200 & $\sim 2\times$ \\
$\zeta$ non-vanishing on $\mathrm{Re}(s) \geq 1$ & 600 & $\sim 2\times$ \\
Dirichlet $L$-functions & 1500 & $\sim 2\times$ \\
Dirichlet's theorem & 800 & $\sim 2\times$ \\
\midrule
Total & $\sim$7400 lines & $\sim 2$--$3\times$ overhead \\
\bottomrule
\end{tabular}
\end{center}

\medskip

\begin{remark}[Caveat about line-count estimates]
The right-hand column is highly speculative; we present it as a
\emph{complexity factor} rather than absolute counts. The factor reflects
the rough overhead of HoTT-native analysis (no classical Mathlib measure
theory, no built-in Lebesgue integration, manual constructive analysis)
plus an honest correction for the immaturity of HoTT analysis libraries.
Actual numbers will diverge from these estimates by easily $\pm 50\%$ as
HoTT analysis libraries mature. The estimates are calibrated against
\cite{Booij2020} for HoTT-native analysis benchmarks.
\end{remark}

%---------------------------------------------------------------
\section{Six Sub-Problems for HoTT-Native \texorpdfstring{$\zeta$}{Zeta}}\label{sec:roadmap}

We propose six concrete sub-problems whose collective resolution would yield
a HoTT-native proof of $\zeta(2) = \pi^2/6$ and a formal HoTT statement of
the Riemann hypothesis.

\subsection*{Sub-problem 1: HoTT-native \texorpdfstring{$\Cc$}{Cc} with full algebraic-closure axiom}
\addcontentsline{toc}{subsection}{Sub-problem 1: HoTT-native complex numbers}
\begin{problem}\label{prob:1}
Construct a HoTT type $\Cc:\mathcal{U}$ that is provably an algebraic closure
of $\Rc$, with the universal property of \Cref{def:UAC}, and verify that
it admits the standard topology and metric structure.
\end{problem}
\textbf{Status:} partially done in HoTT Book setting; algebraic closure not
yet done.

\subsection*{Sub-problem 2: HoTT-native holomorphic functions}
\addcontentsline{toc}{subsection}{Sub-problem 2: HoTT-native holomorphic functions}
\begin{problem}\label{prob:2}
Define \Cref{def:holo} formally; prove the Cauchy-Riemann equations, the
Cauchy integral formula, and the identity theorem (\Cref{thm:Identity}).
\end{problem}
\textbf{Status:} possible by direct constructive analysis. Estimated effort:
$\sim$5000 lines of Agda or Cubical Agda.

\subsection*{Sub-problem 3: HoTT-native Dirichlet series machinery}
\addcontentsline{toc}{subsection}{Sub-problem 3: HoTT-native Dirichlet series}
\begin{problem}\label{prob:3}
Formalize \Cref{def:Dirichlet}, prove \Cref{prop:DirichletAlgebra}, and
construct the basic operations: shift, multiplication, logarithm, derivative.
\end{problem}
\textbf{Status:} our companion Haskell code provides a finite-precision
prototype; HoTT formalisation pending.

\subsection*{Sub-problem 4: HoTT-native analytic continuation}
\addcontentsline{toc}{subsection}{Sub-problem 4: HoTT-native analytic continuation}
\begin{problem}\label{prob:4}
Formalize the analytic-continuation theorem: a holomorphic function on a
connected open admitting a power-series at one boundary point extends
holomorphically to a neighbourhood of that point. Apply to $\zeta$ to obtain
\Cref{def:ZetaUP}.
\end{problem}
\textbf{Status:} this is the key technical bottleneck.

\subsection*{Sub-problem 5: HoTT-native functional equation}
\addcontentsline{toc}{subsection}{Sub-problem 5: HoTT-native functional equation}
\begin{problem}\label{prob:5}
Prove the functional equation \Cref{thm:FE} in HoTT, using either
(a) the Mellin-transform / theta-function method, or (b) Riemann's contour
integral method, or (c) a synthetic cohesive-HoTT proof using analytic-stack
duality.
\end{problem}
\textbf{Status:} requires \Cref{prob:1}--\Cref{prob:4} first.

\subsection*{Sub-problem 6: HoTT-native formal RH statement}
\addcontentsline{toc}{subsection}{Sub-problem 6: HoTT-native RH statement}
\begin{problem}\label{prob:6}
Write down a HoTT proposition $\mathrm{RH}:\Prop$ such that $\mathrm{RH}$
inhabits if and only if every non-trivial zero of $\zeta$ has real part
$1/2$. Verify this is the same statement as the classical RH, modulo the
ambient model.
\end{problem}
\textbf{Status:} doable now (\Cref{conj:RH}), modulo \Cref{prob:1}--\Cref{prob:4}.

\subsection{Composition diagram}

\begin{figure}[htbp]
\centering
\begin{tikzcd}[column sep=small, row sep=large]
\text{Sub-prob 1} \arrow[r] \arrow[dr] & \text{Sub-prob 2} \arrow[d] & \\
& \text{Sub-prob 3} \arrow[r] & \text{Sub-prob 4} \arrow[d] \\
& & \text{Sub-prob 5} \arrow[d] \\
& & \text{Sub-prob 6}
\end{tikzcd}
\caption{Dependency graph for the six sub-problems.}
\label{fig:roadmap}
\end{figure}

\subsection{Estimated effort and milestones}

\begin{itemize}
\item \Cref{prob:1}: $\sim$1000 lines, 1--2 graduate-student-years.
\item \Cref{prob:2}: $\sim$5000 lines, 2--3 g-s-y.
\item \Cref{prob:3}: $\sim$1500 lines, 1 g-s-y.
\item \Cref{prob:4}: $\sim$3000 lines, 2--3 g-s-y; key conceptual work.
\item \Cref{prob:5}: $\sim$5000 lines, 2 g-s-y.
\item \Cref{prob:6}: $\sim$200 lines, weeks. The hard part is the
infrastructure, not the statement.
\end{itemize}

Total: $\sim$15000 lines, 8--12 graduate-student-years. Comparable to the
total effort behind Loeffler--Stoll plus its Mathlib dependencies.

%---------------------------------------------------------------
\section{Open Conjectures: Riemann Hypothesis as a HoTT Statement}\label{sec:rh}

\subsection{Modal-logical status of RH in HoTT}

\Cref{conj:RH} states $\Pi(s:\Cc).\, P(s) \to Q(s)$, where $P$ and $Q$ are
propositions. By the propositional structure, $\mathrm{RH}$ is itself a
proposition (in HoTT-Set sense, an $(-1)$-truncated type).

In particular, $\mathrm{RH}$ is either inhabited or empty; we cannot have
``two non-equivalent proofs'' of RH. This is in contrast to a structural
statement like ``the type of complex numbers admits an algebraic structure''
which has no such uniqueness.

\subsection{Decidability and constructivity}

\begin{remark}\label{rem:decidability}
$\mathrm{RH}$ is \emph{not} decidable in HoTT: we have no algorithm
producing an inhabitant or refutation. This is consistent with
constructivity, since $\mathrm{RH}$ is a $\Pi$-statement over an
uncountable type.
\end{remark}

\begin{remark}\label{rem:LEM}
Under classical logic ($\mathrm{LEM}$), $\mathrm{RH} \vee \neg\,\mathrm{RH}$
is inhabited; this is a consequence of $\mathrm{LEM}$, not a constructive
theorem.
\end{remark}

\subsection{Connections to the Langlands programme}

The Riemann hypothesis is a special case of the Generalized Riemann
Hypothesis (GRH): all non-trivial zeros of all $L$-functions $L(s,\chi)$
attached to primitive Dirichlet characters $\chi$ lie on
$\mathrm{Re}(s) = 1/2$. The Langlands programme conjecturally extends this
to automorphic $L$-functions $L(s,\pi)$.

\begin{conjecture}[GRH for automorphic $L$-functions, HoTT statement]\label{conj:GRH}
\[
\Pi(\pi).\, \Pi(s:\Cc).\, L(s,\pi) = 0 \to
\neg\,\mathsf{trivialZero}_\pi(s) \to \mathrm{Re}(s) = 1/2.
\]
\end{conjecture}

In a HoTT framework where automorphic representations are types in a
suitable $(\infty,1)$-topos, \Cref{conj:GRH} is a $\Pi\Pi\Pi$-statement —
syntactically clean, semantically deep.

\subsection{The von Koch theorem in HoTT}

\begin{theorem}[von Koch, expected HoTT version]\label{thm:vonKoch}
Conditional on \Cref{conj:RH}, the prime-counting function $\pi(x)$
satisfies
\[
|\pi(x) - \mathrm{Li}(x)| \in O(\sqrt{x} \log x).
\]
\end{theorem}

This connects the analytic statement $\mathrm{RH}$ to a discrete one. A
HoTT-native proof would proceed via the explicit formula for $\pi(x)$ in
terms of zeros of $\zeta$, all of which lives downstream of
\Cref{prob:1}--\Cref{prob:6}.

\subsection{What if RH is independent of HoTT-set theory?}

A speculative possibility: $\mathrm{RH}$ might be independent of (some
extension of) HoTT, similar to the way the continuum hypothesis is
independent of ZFC. In that case, HoTT might admit two consistent
extensions: $\mathrm{HoTT} + \mathrm{RH}$ and $\mathrm{HoTT} + \neg\,\mathrm{RH}$.
This is purely speculative and reflects no current consensus; it is offered
only as a contrast with the classical view that RH ``simply has a truth value''.

\subsection{Density theorems and zero-free regions}\label{subsec:zerofree}

Classical analytic number theory has many results about the distribution of
zeros of $\zeta$ short of RH:
\begin{itemize}
\item \textbf{Hadamard / de la Vall\'ee-Poussin (1896):} $\zeta(s) \neq 0$
on $\mathrm{Re}(s) = 1$. Formalized in Lean by Loeffler--Stoll.
\item \textbf{Vinogradov-Korobov:} explicit zero-free region
\[
\mathrm{Re}(s) > 1 - \frac{C}{(\log|t|)^{2/3}\,(\log\log|t|)^{1/3}}.
\]
\item \textbf{Selberg's density theorem:} the proportion of zeros within
distance $\delta$ of the critical line is $\geq 1 - O(\delta^{-1})$.
\item \textbf{Levinson--Conrey:} at least $40\%$ of zeros are on the
critical line.
\end{itemize}

Each of these is, in principle, a HoTT-native theorem once we have
\Cref{prob:1}--\Cref{prob:5}. In particular, the $\mathrm{Re}(s) = 1$
non-vanishing requires only the Euler product and basic
inequalities — no analytic continuation. This is the natural \emph{second}
target after \Cref{thm:Basel}.

\subsection{Effective vs.\ ineffective HoTT statements}\label{subsec:effective}

A subtle point about HoTT: \emph{some} statements about $\zeta$ are
naturally effective (algorithmic), while others are not.

\begin{example}[Effective: $\zeta$ on the line $\mathrm{Re}(s) > 1$]
The map $s \mapsto \sum_{n \leq N} n^{-s}$ for large $N$ provides an
$\epsilon$-approximation to $\zeta(s)$ with explicit $N(\epsilon, s)$
bounds. This is HoTT-native, computable, and compatible with $\Rc$ Cauchy
moduli.
\end{example}

\begin{example}[Ineffective: $\zeta$ in the critical strip]
For $0 < \mathrm{Re}(s) < 1$, the Dirichlet series diverges. Computing
$\zeta(s)$ requires the Riemann--Siegel formula or similar, which involves
contour integration. HoTT-native version would require \Cref{prob:4}.
\end{example}

\begin{example}[Ineffective: trivial zero locations]
The trivial zeros at $s = -2, -4, -6, \ldots$ are \emph{deduced} from the
functional equation; their existence requires \Cref{prob:5}.
\end{example}

These distinctions matter for HoTT formalization: parts of the theory are
algorithmic-friendly (Cubical Agda can compute them), other parts are not.

%---------------------------------------------------------------
\section{Discussion}\label{sec:discussion}

\subsection{Why is analytic NT specifically hard?}

Algebraic objects (rings, modules, groups, schemes) are \emph{small}: their
data fit on a finite page or in a computer's memory. Analytic objects
($\R$, $\C$, holomorphic functions, contour integrals) are \emph{large} in
two senses:
\begin{itemize}
\item Type-theoretic size: $\Rc$ is uncountable, so any property of $\Rc$
is in some sense a $\Pi$-statement over an uncountable type.
\item Logical complexity: $\zeta$ involves nested quantifiers over reals,
limits of limits, integrals, and analytic continuations.
\end{itemize}
HoTT, like any type theory, is well-suited to data of small or medium
complexity; encoding heavy analysis requires either substantial primitive
infrastructure or a transport from a classical foundation.

\subsection{Comparison: Lean Mathlib's effectiveness}

Lean Mathlib's success at formalizing analytic NT (Loeffler--Stoll 2025)
relies on:
\begin{enumerate}
\item Classical logic everywhere.
\item A large prebuilt library of measure theory, integration, and complex
analysis.
\item Decidability assumptions where appropriate.
\item Set-theoretic ambient foundations (Lean 4 type theory is
ZFC-equivalent in strength).
\end{enumerate}
HoTT can match (1)--(3) at extra cost; (4) is a deeper foundational issue.

\subsection{The role of cohesive / differential HoTT}

\emph{Differential cohesive HoTT}~\cite{Schreiber2013, Wellen2018} extends
HoTT with shape, flat, and sharp modalities. In this setting, smooth /
holomorphic / meromorphic become synthetic. Examples like
\cite{ScreiberShulman2014QFT} show that synthetic differential geometry is
viable in HoTT.

A natural research direction is to lift the Loeffler--Stoll work into
\emph{differential cohesive HoTT}, treating $\C$ as a $\esh$-modal type and
$\zeta$ as a synthetic meromorphic function. This is a substantial project
but appears feasible.

\subsection{Limitations of this paper}

\begin{itemize}
\item We do not give a working HoTT formalization; we give a research
roadmap and partial Haskell prototypes.
\item We do not prove the conjectures of \Cref{thm:ZetaEquiv} or
\Cref{thm:FE}; we sketch proofs that depend on currently-unavailable
HoTT-native infrastructure.
\item Our proposed bridge (\Cref{prob:bridge}) between condensed mathematics
and HoTT is open; we have not constructed the bridging $(\infty,1)$-topos.
\end{itemize}

\subsection{Future directions}

\begin{enumerate}
\item Implement \Cref{prob:1}--\Cref{prob:3} in Cubical Agda.
\item Develop the cohesive-HoTT bridge for analytic continuation
(\Cref{prob:4}).
\item Translate Loeffler--Stoll Lean to HoTT-axiomatised classical $\C$ as
a first benchmark.
\item Use the synthesis from~\cite{Long2026Synthesis} to attack
\Cref{conj:GRH} via geometric Langlands in $(\infty,1)$-topoi.
\item Explore the analytic Langlands programme of EFK
\cite{EFK2025AnalyticLanglands} in HoTT.
\end{enumerate}

\subsection{Connection with directed type theory}

A recent development is \emph{directed type theory} (DTT) of Riehl,
Shulman, Verity, North~\cite{Riehl2021Elements}. DTT extends HoTT with a
notion of directed morphism — paths that are not invertible. This gives a
potential synthetic foundation for $(\infty,1)$-categories.

Why is DTT relevant to analytic NT?
\begin{itemize}
\item Automorphic representations are functors, not equivalences. DTT
captures functorial structure natively.
\item Hecke operators on $L$-functions act as endomorphisms, not
automorphisms; DTT distinguishes these.
\item The Langlands functor $\sigma:\mathsf{Aut}(GL_n) \to \mathsf{Galois}$
is a functor between $\infty$-categories of representations; DTT gives a
type-theoretic home.
\end{itemize}

\subsection{Connection with synthetic algebraic geometry}

\emph{Synthetic algebraic geometry}~\cite{CherubiniRijke2024} works inside
HoTT internally to a Zariski (or étale) topos, and develops algebraic
geometry without external schemes. The relevant facts:
\begin{itemize}
\item Schemes become types satisfying a Zariski-locality condition.
\item Sheaf cohomology is internal cohomology.
\item Group schemes, including reductive groups for Langlands, become
types in the internal language.
\end{itemize}

A natural research line: use synthetic AG \emph{plus} cohesion to define
analytic moduli stacks, and develop $\zeta$ as an internal object.

\subsection{The role of computer-checked proofs}

A practical question: \emph{should} we even attempt HoTT formalisation of
analytic NT, given that classical Lean Mathlib is so much more advanced?

Three answers:
\begin{enumerate}
\item \emph{Foundational interest:} HoTT and Lean differ at the foundation
level. Univalence simplifies certain isomorphism / equivalence arguments.
A HoTT formalisation would be qualitatively different.
\item \emph{Internal-language gain:} working inside an $(\infty,1)$-topos
gives access to the topos's internal logic. For instance, condensed
mathematics has its own internal language (the Solid topos), distinct from
classical set theory; HoTT might bridge these.
\item \emph{Educational value:} a HoTT formalisation forces clarity about
what data $\zeta$ ``really is''. The exercise of distinguishing the three
candidate definitions and identifying their pairwise equivalences is
foundationally illuminating.
\end{enumerate}

%---------------------------------------------------------------
\section{Conclusion}\label{sec:conclusion}

We have formulated, with technical care, the principal open problem of our
prior synthesis: HoTT-native analytic number theory. The Riemann zeta
function does not yet exist as a HoTT-native object; the statement
$\zeta(s)=0$ is not yet a HoTT-native proposition. Closing this gap
requires the prerequisite chain of \Cref{sec:prereq}, the candidate
definitions of \Cref{sec:zeta}, the $(\infty,1)$-topos perspective of
\Cref{sec:langlands}--\Cref{sec:condensed}, the comparison with
Loeffler--Stoll Lean of \Cref{sec:loefflerstoll}, and the six concrete
sub-problems of \Cref{sec:roadmap}.

We do not solve the principal open problem. We formulate it precisely, show
it is not vacuous, and provide enough structure for future workers to take
concrete steps. The roadmap of \Cref{sec:roadmap} estimates 8--12
graduate-student-years to deliver a HoTT-native analogue of Loeffler--Stoll
2025. This is significant, but tractable; the much harder Riemann
hypothesis itself remains untouched.

The dialogue between HoTT and analytic number theory is just beginning.
Geometric Langlands has been proven (2024) using $(\infty,1)$-categorical
machinery that overlaps with HoTT's intended models. Condensed mathematics
provides a uniform setting for analytic objects. These two streams will,
we conjecture, converge on a HoTT-native analytic number theory in the
coming decade. This paper aims to map the territory.

%---------------------------------------------------------------
\appendix

\section{Appendix: HoTT axioms for analytic NT}\label{app:axioms}

For reference, we collect the HoTT axioms / definitions that any
HoTT-native analytic NT should respect. We work in MLTT with one
universe $\mathcal{U}$ and the following axioms.

\subsection*{A.1 Univalence}

\begin{equation}\label{eq:univ}
\ua : \Pi(A,B:\mathcal{U}).\, (A \simeq B) \simeq (A = B).
\end{equation}

This is the standard univalence axiom of~\cite{HoTTBook}. It implies
function extensionality and propositional extensionality.

\subsection*{A.2 Higher inductive types}

We assume HITs of the form needed in this paper:
\begin{itemize}
\item Cauchy reals $\Rc$ as a HIIT (\Cref{def:Rc}).
\item Set quotients $A/{\sim}$ for relations $\sim$ on $A$.
\item Propositional truncation $\|A\|_{-1}$.
\item $n$-truncations $\|A\|_n$.
\end{itemize}

\subsection*{A.3 Cohesion (optional)}

For synthetic-analytic statements, cohesive HoTT
\cite{Shulman2018Cohesive} adds:
\begin{itemize}
\item Shape modality $\esh:\mathcal{U} \to \mathcal{U}$ (left adjoint to
$\flat$).
\item Flat modality $\flat:\mathcal{U} \to \mathcal{U}$ (right adjoint to
$\esh$).
\item Sharp modality $\sharp:\mathcal{U} \to \mathcal{U}$.
\end{itemize}

These satisfy the cohesive triple-adjunction $\esh \dashv \flat \dashv \sharp$
and natural-transformation laws between them.

\subsection*{A.4 Choice (optional, classical)}

For some classical analytic-NT results we may need:
\begin{equation}\label{eq:choice}
\mathrm{AC} : \Pi(A:\mathcal{U}).\, \Pi(B:A \to \mathcal{U}).\,
   (\Pi(a:A).\, \|B(a)\|_{-1}) \to \|\Pi(a:A).\, B(a)\|_{-1}.
\end{equation}

Equivalent to set-theoretic AC for sets. Required for classical existence
of algebraic closures, classical functional equation proofs.

\subsection*{A.5 Excluded middle (optional, classical)}

\begin{equation}\label{eq:lem}
\mathrm{LEM} : \Pi(P:\Prop).\, P + \neg P.
\end{equation}

Required for the statement of $\mathrm{RH} \vee \neg\,\mathrm{RH}$ as
``there is a fact of the matter''.

\section{Appendix: Companion code overview}\label{app:code}

The companion repository at \texttt{src/langlands-analytic-number-theory/}
contains:

\begin{itemize}
\item \texttt{Main.hs} --- driver illustrating partial $\zeta$ sums and
Dirichlet series operations.
\item \texttt{Zeta.hs} --- partial zeta computations with explicit
truncation.
\item \texttt{Dirichlet.hs} --- Dirichlet series machinery, Dirichlet
convolution, multiplicative functions.
\item \texttt{Properties.hs} --- QuickCheck properties for the algebraic
laws on Dirichlet series.
\item \texttt{Proofs.hs} --- equational proofs of the Euler product
identity, modulo a finite-truncation axiom.
\end{itemize}

The Lean component at \texttt{lean/langlands-analytic-number-theory/}
contains:

\begin{itemize}
\item \texttt{Zeta.lean} --- a Lean 4 / Mathlib sketch indexed against
Loeffler--Stoll, marking which lemmas would translate cleanly to HoTT.
\end{itemize}

These are illustrative prototypes, not formal verifications. Their purpose
is to make the prerequisite chain (\Cref{sec:prereq}) and the candidate
definitions (\Cref{sec:zeta}) executable, so that future HoTT-native work
has computational benchmarks.

\section{Appendix: Worked sub-roadmap for sub-problem 4}\label{app:cont}

We sketch a plausible HoTT proof strategy for analytic continuation of
$\zeta$, addressing \Cref{prob:4} in more detail.

\subsection*{C.1 The Hurwitz integral representation}

Classical analytic continuation of $\zeta$ via the Hurwitz integral
representation~\cite{Hurwitz1882}:
\[
\zeta(s) = \frac{1}{\Gamma(s)} \int_0^\infty \frac{x^{s-1}}{e^x - 1}\, dx,
\qquad \mathrm{Re}(s) > 1.
\]
The integrand is integrable for $\mathrm{Re}(s) > 1$ near $0$ and is
exponentially small at $\infty$. To analytically continue, one deforms the
contour:
\[
\zeta(s) = \frac{\Gamma(1-s)}{2\pi i}\,\oint_C \frac{(-x)^{s-1}}{e^x - 1}\,dx,
\]
where $C$ is the Hankel contour: from $+\infty$ along $\R^+$, around $0$
counterclockwise, back to $+\infty$.

\subsection*{C.2 What HoTT needs}

To carry out the contour-integral proof in HoTT we need:
\begin{enumerate}
\item HoTT-native $\Gamma$ function: $\Gamma(s) := \int_0^\infty t^{s-1}
e^{-t}\, dt$ extended via the functional equation $\Gamma(s+1) = s\Gamma(s)$
to all $\Cc \setminus \{0,-1,-2,\ldots\}$.
\item HoTT-native contour integration: a path $\gamma:[0,1] \to \Cc$ and an
integral $\int_\gamma f(z)\, dz$, satisfying linearity, change-of-variables,
and Cauchy's theorem.
\item HoTT-native version of Cauchy's theorem: holomorphic functions integrate
to zero around closed contours.
\item HoTT-native deformation lemma: integral over homotopic contours agree.
\item HoTT-native passage from real integral to contour integral
(Mellin-Barnes type).
\end{enumerate}

\subsection*{C.3 Estimated effort}

Each item above is a self-contained HoTT formalisation problem of moderate
size. Cumulatively, this is the heart of \Cref{prob:4}, and we estimate
it at $\sim 3000$ lines of Cubical Agda.

\subsection*{C.4 Alternative strategies}

\begin{itemize}
\item \emph{Riemann's theta-function method}: $\zeta(s)$ via Mellin transform
of $\theta(t) := \sum_{n \in \Z} e^{-\pi n^2 t}$. Requires the modular
transformation $\theta(t) = t^{-1/2}\theta(1/t)$.
\item \emph{Hardy's symmetric form}: $\xi(s) := \frac{s(s-1)}{2}\pi^{-s/2}
\Gamma(s/2)\zeta(s)$, satisfying $\xi(s) = \xi(1-s)$. Cleaner but uses the
same infrastructure.
\item \emph{Cohesive synthetic approach}: Inside differential cohesive HoTT
\cite{Schreiber2013}, define $\zeta$ via the synthetic functional integral.
Speculative.
\end{itemize}

\section{Appendix: A glossary of HoTT-native synonyms}\label{app:glossary}

For readers more familiar with classical analytic NT than with HoTT, we
provide a short translation glossary:

\medskip

\begin{center}
\begin{tabular}{ll}
\toprule
Classical & HoTT-native \\
\midrule
$\R$ & $\Rc$ (Cauchy reals) \\
$\C$ & $\Cc$ (univalent algebraic closure of $\Rc$) \\
$f:\C \to \C$ holomorphic & $f$ admits constructive Newton quotient \\
``unique up to isomorphism'' & propositionally equal under univalence \\
``the algebraic closure'' & ``a univalent algebraic closure'' \\
ZFC's choice & propositional choice axiom (optional) \\
LEM & $\Pi(P:\Prop).\, P + \neg P$ (optional) \\
quotient set & set quotient $A/{\sim}$ via HIT \\
limit of sequences & limit constructor of HIIT $\Rc$ \\
``open subset'' & ``cohesive open'' or ``(-1)-truncated open'' \\
``connected'' & propositional truncation of $A$ is contractible \\
sheaf cohomology & internal cohomology of cohesive HoTT \\
$\infty$-category & type with Segal condition \\
$\infty$-topos & univalent universe in cohesive HoTT \\
``functor'' & morphism in directed type theory \\
``natural transformation'' & 2-cell in directed type theory \\
condensed set & sheaf on profinite types in cohesive HoTT (open) \\
analytic stack & 6-functor formalism in cohesive HoTT (open) \\
$\zeta$ classically & no HoTT-native definition yet \\
\bottomrule
\end{tabular}
\end{center}

\medskip

The bottom row is the principal open problem of this paper.

%---------------------------------------------------------------
\bibliographystyle{plain}
\begin{thebibliography}{99}

\bibitem{HoTTBook}
Univalent Foundations Program.
\textit{Homotopy Type Theory: Univalent Foundations of Mathematics}.
Institute for Advanced Study, 2013. \url{https://homotopytypetheory.org/book}

\bibitem{Long2026Synthesis}
M.\ Long and the YonedaAI Collaboration.
\textit{Synthesis: Six Open Problems at the Interface of HoTT and
Twenty-First Century Mathematics}.
YonedaAI Research Collective, 2026.

\bibitem{Voevodsky2014}
V.\ Voevodsky.
\textit{Univalent foundations of mathematics}.
ICM Talk, 2014.

\bibitem{Booij2020}
A.\ Booij.
\textit{Analysis in univalent type theory}.
PhD thesis, University of Birmingham, 2020.

\bibitem{Mines1988}
R.\ Mines, F.\ Richman, W.\ Ruitenburg.
\textit{A Course in Constructive Algebra}.
Springer Universitext, 1988.

\bibitem{Shulman2018Cohesive}
M.\ Shulman.
\textit{Brouwer's fixed-point theorem in real-cohesive homotopy type theory}.
Mathematical Structures in Computer Science, 28:856--941, 2018.

\bibitem{Shulman2019InfTopos}
M.\ Shulman.
\textit{All $(\infty,1)$-toposes have strict univalent universes}.
arXiv:1904.07004, 2019.

\bibitem{Lurie2009HTT}
J.\ Lurie.
\textit{Higher Topos Theory}.
Annals of Mathematics Studies 170, Princeton University Press, 2009.

\bibitem{LurieHA}
J.\ Lurie.
\textit{Higher Algebra}.
Manuscript, 2017. \url{https://www.math.ias.edu/~lurie/papers/HA.pdf}

\bibitem{Gaitsgory2024GLC}
D.\ Gaitsgory et al.
\textit{Proof of the geometric Langlands conjecture I--V}.
arXiv:2405.03599 and sequels, 2024.
\url{https://www.mpim-bonn.mpg.de/gaitsgde/GLC/}

\bibitem{BeilinsonDrinfeld2004}
A.\ Beilinson, V.\ Drinfeld.
\textit{Chiral Algebras}.
American Mathematical Society Colloquium Publications 51, 2004.

\bibitem{LoefflerStoll2025}
D.\ Loeffler, M.\ Stoll.
\textit{Formalizing zeta and L-functions in Lean}.
arXiv:2503.00959. Annals of Formalized Mathematics, vol.\ 1, 2025.

\bibitem{Scholze2019Condensed}
P.\ Scholze.
\textit{Lectures on condensed mathematics}.
Lecture notes, University of Bonn, 2019.
\url{https://www.math.uni-bonn.de/people/scholze/Condensed.pdf}

\bibitem{ClausenScholze2024Analytic}
D.\ Clausen, P.\ Scholze.
\textit{Analytic stacks}.
Lecture series, MPIM Bonn, 2023--2024.

\bibitem{ClausenScholze2024SixFunctor}
D.\ Clausen, P.\ Scholze.
\textit{Six-functor formalisms in analytic geometry}.
arXiv:2507.08566, 2025.

\bibitem{ClausenScholze2025AnalyticStacks}
D.\ Clausen, P.\ Scholze.
\textit{Analytic stacks (continued)}.
arXiv:2512.14612, 2025.

\bibitem{ClausenScholzeComplex}
D.\ Clausen, P.\ Scholze.
\textit{Condensed mathematics and complex geometry}.
\url{https://people.mpim-bonn.mpg.de/scholze/Complex.pdf}, 2022.

\bibitem{EFK2025AnalyticLanglands}
P.\ Etingof, E.\ Frenkel, D.\ Kazhdan.
\textit{An analytic version of the Langlands correspondence over the
function field}.
arXiv:2103.01509 and follow-ups, 2021--2025.

\bibitem{CisinskiKKNS2025}
D.-C.\ Cisinski et al.
\textit{Type-theoretic foundations of $(\infty,1)$-categories}.
Manuscript / arXiv preprints, 2025.

\bibitem{Langlands1970}
R.\ P.\ Langlands.
\textit{Problems in the theory of automorphic forms}.
Lecture Notes in Mathematics 170, Springer, 1970.

\bibitem{GelbartLanglands1979}
S.\ Gelbart.
\textit{An elementary introduction to the Langlands program}.
Bull.\ AMS, 10(2):177--219, 1984.

\bibitem{Schreiber2013}
U.\ Schreiber.
\textit{Differential cohomology in a cohesive infinity-topos}.
arXiv:1310.7930, 2013.

\bibitem{Wellen2018}
F.\ Wellen.
\textit{Cartan geometry in modal homotopy type theory}.
PhD thesis, Karlsruhe Institute of Technology, 2018.

\bibitem{ScreiberShulman2014QFT}
U.\ Schreiber, M.\ Shulman.
\textit{Quantum gauge field theory in cohesive homotopy type theory}.
EPTCS 158, 2014.

\bibitem{Euler1737}
L.\ Euler.
\textit{Variae observationes circa series infinitas}.
Commentarii academiae scientiarum Petropolitanae 9, 1737.

\bibitem{RiemannHypothesisOriginal}
B.\ Riemann.
\textit{Über die Anzahl der Primzahlen unter einer gegebenen Größe}.
Monatsberichte der Königlich Preußischen Akademie der Wissenschaften zu
Berlin, 1859.

\bibitem{KapustinWitten2007}
A.\ Kapustin, E.\ Witten.
\textit{Electric-magnetic duality and the geometric Langlands program}.
Communications in Number Theory and Physics, 1(1):1--236, 2007.

\bibitem{Riehl2021Elements}
E.\ Riehl, D.\ Verity.
\textit{Elements of $\infty$-Category Theory}.
Cambridge Studies in Advanced Mathematics 194, Cambridge University Press, 2022.

\bibitem{CherubiniRijke2024}
F.\ Cherubini, E.\ Rijke et al.
\textit{Synthetic algebraic geometry in homotopy type theory}.
Manuscript, \url{https://github.com/felixwellen/synthetic-zariski}, 2024.

\bibitem{Hurwitz1882}
A.\ Hurwitz.
\textit{Einige Eigenschaften der Dirichlet'schen Funktionen $F(s)$}.
Zeitschrift für Mathematik und Physik, 27:86--101, 1882.

\end{thebibliography}

\end{document}
